The implementation costs of algorithmic redistricting using the \texttt{bisect} platform are a minor fraction of any alternative. At \$50,000--\$75,000 per cycle---versus \$3 million for the Iowa model and \$35 million for the California model---algorithmic redistricting is not merely cheaper but an order of magnitude cheaper. The comparison to redistricting litigation costs (\$10--\$40 million in contested states) makes the cost-benefit case overwhelming: the cost of implementing algorithmic redistricting is less than the legal fees for a single redistricting motion in major litigation.

The dominant cost driver is not the algorithm---compute costs are under \$20---but staff time for quality assurance, public engagement, and legal review. These costs are real and should not be minimized: algorithmic redistricting requires competent human oversight to verify outputs, document departures, and ensure VRA compliance. But these costs are substantially lower than in any alternative process because the algorithm eliminates the months-long map-drawing phase that generates the largest staff expenditures in traditional redistricting.

The \texttt{bisect} manifest format satisfies APA record-keeping requirements directly. The manifest provides complete data provenance (which census release and TIGER file were used), documents criteria application (which parameters implement which criteria), and enables departure documentation (the Commission Adjustment Protocol from P.3). The 12-year data retention requirement is achievable at negligible cost (\$20 in cloud storage) and satisfies the period during which redistricting decisions can be challenged.

The model procurement specification in Section~5 gives state redistricting offices a ready-to-adapt document for contracting algorithmic redistricting services. It establishes minimum technical requirements (open-source or escrowed code; reproducible manifest; partisan-blind algorithm) that protect states against vendors whose tools are technically unsound or manipulable. The evaluation criteria balance technical quality, price, and past performance in a way that favors well-documented, verifiable approaches.

For administrators and budget officers reading this paper: the question is not whether your state can afford algorithmic redistricting. The question is whether your state can afford the alternative. States that have spent double-digit millions on redistricting litigation could have funded algorithmic redistricting for 200 cycles. The budget case for algorithmic redistricting is as strong as the legal and political cases developed in P.0 through P.4. Together, the five case papers in Track P provide a complete implementation toolkit: the political economy analysis (P.0), the federal legislative template (P.1), the ballot initiative design (P.2), the commission integration protocol (P.3), the court order models (P.4), and the cost accounting and procurement specification (P.5). The toolkit is ready. The political conditions for using it vary by state and by moment. But the tool itself is complete, operational, and affordable.
